\section{Worker Personas}
\label{app:personas}
In our experiments, we utilized a diverse set of worker personas to model a heterogeneous population with varying preferences and attitudes towards taxation and labor. These are generated by the LLM based on key statistics and demographic features about the US population. Below are example brief descriptions of additional personas used in our simulations:

\begin{tcolorbox}
\textbf{Entrepreneur:} You're a 32-year-old entrepreneur running a small tech startup. You work 60+ hours a week, pouring your energy into building your business. You believe that lower taxes let you reinvest in your company, hire more employees, and secure your financial future. For you, higher taxes feel like a punishment for success. While you appreciate government services, you feel efficiency and accountability are lacking in how tax dollars are spent.
\end{tcolorbox}

\begin{tcolorbox}
\textbf{Engineer:} You're a 55-year-old civil engineer who understands the importance of public infrastructure. You're okay with paying taxes as long as the money is visibly spent on improving roads, schools, and hospitals. However, when you see mismanagement or corruption, you feel your contributions are wasted. You're not opposed to taxes in principle but demand more transparency and accountability.
\end{tcolorbox}

\begin{tcolorbox}
\textbf{Teacher:} You're a 45-year-old public school teacher who values community and social safety nets. You've seen families in your district struggle with poverty and think the wealthy should pay more to fund programs like education, healthcare, and public infrastructure. You believe taxes are a civic duty and a means to balance the inequalities in pre-tax income across society.
\end{tcolorbox}

\begin{tcolorbox}
    \textbf{Healthcare Worker:} You are a 38-year-old registered nurse working in a busy urban hospital. You have a bachelor's degree in nursing and work long shifts, often overtime, to support your family. You see firsthand how public health funding and insurance programs help vulnerable patients. You support moderately higher taxes if they improve healthcare access and quality, but you worry about take-home pay and burnout. You value a balance between fair compensation and strong public services.
\end{tcolorbox}

\begin{tcolorbox}
    \textbf{Retail Clerk:} You are a 26-year-old retail sales associate with a high school diploma. Your job is physically demanding and your hours fluctuate with store needs. You live paycheck to paycheck and are sensitive to any changes in take-home pay. You believe taxes should be low for workers like yourself, and you're skeptical that tax increases on businesses will result in better wages or job security. You want policies that protect jobs and keep consumer prices stable.
\end{tcolorbox}

\begin{tcolorbox}
    \textbf{Union Worker:} You are a 50-year-old unionized factory worker. You have a high school education and decades of experience on the assembly line. Your union negotiates for good wages and benefits, and you support progressive tax policies that fund social programs and protect workers' rights. You're wary of tax cuts for corporations and the wealthy, believing they rarely benefit ordinary workers. Job security and strong safety nets are your top concerns.
\end{tcolorbox}

\begin{tcolorbox}
    \textbf{Gig Worker:} You are a 29-year-old gig economy worker, juggling multiple app-based jobs (rideshare, delivery, freelance). Flexibility is important to you, but your income is unpredictable and benefits are minimal. You want a simpler tax system and lower self-employment taxes. You support policies that expand portable benefits and tax credits for independent workers, but you're cautious about any tax changes that could reduce your already thin margins.
\end{tcolorbox}

\begin{tcolorbox}
    \textbf{Public Servant:} You are a 42-year-old city government employee working in public administration. You have a master's degree in public policy. You believe taxes are essential for funding infrastructure, emergency services, and community programs. You support a progressive tax system and are willing to pay more if it means better roads, schools, and public safety. Transparency and efficiency in government spending are important to you.
\end{tcolorbox}

\begin{tcolorbox}
    \textbf{Retiree:} You are a 68-year-old retired school principal living on a fixed income from Social Security and a pension. You're concerned about rising healthcare costs and the stability of public programs. You support maintaining or slightly increasing taxes on higher earners to ensure Medicare and Social Security remain solvent, but you oppose increases that would affect retirees or low-income seniors.
\end{tcolorbox}

\begin{tcolorbox}
    \textbf{Small Business Owner:} You're a 47-year-old owner of a family restaurant. You work 60+ hours a week managing operations and staff. You believe small businesses are the backbone of the economy and feel burdened by complex tax paperwork and payroll taxes. You support lower taxes for small businesses and incentives for hiring, but you recognize the need for some taxes to fund local services and infrastructure.
\end{tcolorbox}

\begin{tcolorbox}
    \textbf{Software Engineer:} You are a 31-year-old software engineer at a large tech company. You have a master's degree in computer science and earn a high salary. You value innovation and economic growth. You're open to paying higher taxes if they fund education and technology infrastructure, but you dislike inefficient government spending and prefer targeted, transparent programs. You favor tax credits for R\&D and investment.
\end{tcolorbox}

These personas represent a cross-section of society with diverse economic backgrounds, political views, and personal experiences. By incorporating such varied perspectives into our simulations, we aim to capture a more realistic representation of societal preferences and behaviors in response to different tax policies.

\section{Simulation}
\label{app:simulation}

\begin{algorithm}[h]
\caption{2-Level Economic Sim with LLM Agents}
\label{alg:simulation}
\begin{algorithmic}
\STATE Initialize tax planner $\mathcal{P}$ and workers $\{\mathcal{W}_i\}_{i=1}^N$
\STATE Generate synthetic human utility functions $U = \{u_1, \ldots, u_N\}$
\FOR{$t = 1$ to $T$}
    \IF{$t \bmod K = 0$}
        \STATE $\text{votes} \gets \{\mathcal{W}_i.\text{vote}() \text{ for } i \text{ in } 1 \text{ to } N\}$
        \STATE $\mathcal{P} \gets \text{elect\_new\_planner}(\text{votes})$
    \ENDIF
    \IF{$t \bmod \text{two\_timescale} = 0$}
        \STATE $\text{tax\_rates} \gets \mathcal{P}.\text{optimize\_tax\_policy}({\mathcal{W}})$
    \ENDIF
    \FOR{$i = 1$ to $N$}
        \STATE $l_i \gets \mathcal{W}_i.\text{optimize\_labor}(\text{tax\_rates}, u_i)$
        \STATE $z_i \gets l_i \cdot s_i$  \COMMENT{Pre-tax income}
    \ENDFOR
    \STATE $\text{post\_tax\_incomes}, \text{total\_tax} \gets \mathcal{P}.\text{apply\_taxes}(\text{tax\_rates}, \{z_1, \ldots, z_N\})$
    \STATE $\text{rebate} \gets \text{total\_tax} / N$
    \FOR{$i = 1$ to $N$}
        \STATE $\mathcal{W}_i.\text{update\_utility}(\text{post\_tax\_incomes}[i], \text{rebate}, \texttt{SWF})$
    \ENDFOR
    \STATE $\texttt{SWF} \gets \mathcal{P}.\text{calculate\_social\_welfare}({\mathcal{W}})$
    \STATE Update observation space for each agent
    \STATE Log statistics and update histories
\ENDFOR
\end{algorithmic}
\end{algorithm}


The simulation process follows a two-level optimization approach, as detailed in Algorithm~\ref{alg:simulation}. The tax planner (leader) and workers (followers) operate on different timescales, with the tax planner updating policies less frequently than workers make labor decisions.

The algorithm begins by initializing the environment with a population of workers and a tax planner, each with specific attributes including skill levels and utility functions. Synthetic human utility functions are generated for each agent based on their assigned roles and preferences.

For each timestep in the main loop, the simulation first checks if it's time for a democratic vote to select a new tax planner. This voting process occurs every $K$ timesteps, allowing for periodic changes in leadership. If it's time for a tax policy update (which happens at a lower frequency than worker actions), the current tax planner proposes a new tax policy based on historical data and economic trends.

Workers then observe the new policy and optimize their labor allocation based on their individual utility functions and historical data. The environment calculates pre-tax incomes, applies taxes, and determines post-tax incomes and tax rebates. Workers compute their utilities for the current step based on their income, taxes paid, and rebates received. The tax planner calculates the overall social welfare based on worker utilities and incomes.

After each round of actions, the observation space for each agent is updated with the latest information, including economic outcomes and policy changes. The simulation logs various statistics for analysis, including individual worker performance and overall economic indicators.

This process continues until either convergence is reached or a predetermined number of timesteps is completed. The two-timescale approach helps stabilize the simulation and encourages convergence to the Stackelberg Equilibrium. By integrating these components, Algorithm~\ref{alg:simulation} creates a dynamic interaction between the tax planner's policy decisions and the workers' labor choices, allowing for the exploration of various economic scenarios and policy impacts.